\documentclass[]{article}

\usepackage{amsmath}
\usepackage{multirow}
\usepackage{natbib}
\usepackage{graphicx} 


\begin{document}



In this work, we addressed the question of whether the anomalous transport of passive tracers in a deterministic, chaotic Hamiltonian system can be fully described by a canonical stochastic model. To this end, we performed a direct statistical comparison between tracer trajectories simulated in a two-dimensional point-vortex flow and trajectories from a benchmark Lévy walk model. The chaoticity of the vortex system was controlled by varying the number of vortices, $N$, while the transport regime of the Lévy walk was set by its power-law tail exponent, $\beta$.

Our analysis employed a suite of statistical measures, including the mean squared displacement, velocity autocorrelation functions, residence time distributions, and displacement probability density functions. The results demonstrate that as the vortex density increases, the tracer transport transitions from near-normal diffusion ($\alpha \approx 1.05$ for $N=5$) to strong superdiffusion ($\alpha \approx 1.52$ for $N=40$). This transition is mechanistically supported by the emergence of power-law distributions for residence times in low-velocity regions and the development of heavy, non-Gaussian tails in the displacement distributions. These kinematic signatures closely mimic those of a superdiffusive Lévy walk, suggesting a strong correspondence at the level of macroscopic transport.

Despite these kinematic similarities, a fundamental divergence was uncovered when we analyzed the ergodicity of the systems. We found that the point-vortex system becomes progressively more ergodic as the superdiffusive behavior strengthens with increasing $N$. In this chaotic flow, a higher degree of chaos leads to more efficient phase-space sampling by individual tracers, reducing the statistical variance across the ensemble. This trend is in direct opposition to the behavior of the Lévy walk model, which becomes increasingly non-ergodic as its superdiffusive character intensifies. For Lévy walks, stronger superdiffusion is driven by rarer and longer ballistic flights, which inherently increases the statistical heterogeneity among individual trajectories.

The central conclusion of this study is that while a chaotic point-vortex flow can reproduce the key kinematic features of a Lévy walk, the two systems are not statistically equivalent. The opposing relationship between the degree of superdiffusion and the degree of ergodicity breaking reveals a fundamental difference rooted in their underlying dynamics: one is a deterministic, Hamiltonian system with memory, while the other is a memoryless stochastic process. This finding underscores that matching macroscopic scaling laws, such as the mean squared displacement exponent, is an insufficient condition for validating a stochastic model against a complex deterministic system. A complete description must also account for higher-order statistical properties that reflect the structure of the underlying dynamical ensemble.

\end{document}
                